\paragraph{World and context modeling.}
The project models the therapy room through a hybrid semantic scene-graph representation, rather than through a single monolithic map or ontology. Following the semantic-map view discussed in the course, the representation can be interpreted as a tuple $\langle R, M, P\rangle$: $R$ is the shared global reference frame, implemented as the ROS/Nav2 \texttt{map} frame; $M$ contains metric information such as navigation poses, 3D object centroids, bounding boxes, and drop poses; and $P$ contains semantic predicates and labels describing what those metric entities mean for the task. In the implementation this is split into two complementary scene graphs. The first one is the manually defined static semantic pose graph in \texttt{src/empower/config/scenes/therapy\_room\_scene\_graph.json}. It stores stable prior knowledge about the therapy setting: the room, the therapist, the child, the room center used as capture origin, the child approach waypoint, and the child delivery table. Nodes have semantic types such as \texttt{room}, \texttt{person}, \texttt{zone}, \texttt{nav\_pose}, \texttt{waypoint}, and \texttt{surface}; edges encode predicates such as \texttt{located\_in}, \texttt{approaches}, \texttt{reachable\_via}, and \texttt{used\_for}. The second graph is the runtime EMPOWER graph, written as \texttt{global\_scene\_graph.json} together with \texttt{global\_detections.json}. This graph is generated from perception and contains the visual objects detected during the current multi-view scan, with labels, aliases, confidence scores, source views, 3D centroids in the map frame, and 3D bounding boxes.

This separation is important for task execution. The static graph represents social and functional priors that are known before perception starts and that are difficult or unnecessary to infer visually in this simulation. For example, the mission orchestrator resolves the therapist interaction pose, the center/capture pose, and the child delivery pose from the static graph, so that TIAGo can first approach the therapist, receive the task, return to the capture origin, and finally reach the child. The low-level executor uses the same graph to ground language-level targets such as ``child'', ``kid'', ``person'', or ``therapist'' into concrete Nav2 goals. It also resolves the intermediate child waypoint and the delivery-table drop pose. In contrast, navigation toward objects such as the book, teddy bear, or computer is grounded through the EMPOWER runtime graph: the executor loads the current global detections, matches symbolic plan targets against object labels and aliases, retrieves the selected object's 3D centroid, and computes a standoff navigation goal. In this way, the planner can use symbolic actions such as \texttt{NAVIGATE}, \texttt{GRAB}, and \texttt{DROP}, while execution remains grounded in metric coordinates.

Knowledge is acquired from multiple sources. Predefined knowledge comes from the Gazebo world, the occupancy map, configuration files, object aliases, and the static scene graph. Online visual knowledge is acquired by the EMPOWER pipeline: TIAGo performs a multi-view RGB-D capture, YOLO-World proposes open-vocabulary object detections, EfficientViT-SAM provides segmentation masks, depth is projected into point clouds, and detections from different views are transformed into the \texttt{map} frame and merged according to semantic label and spatial consistency thresholds. User input contributes a second contextual layer: the therapist provides the child's name, can specify or override the requested activity, and can determine the physical task through speech or text published on \texttt{/hri/user\_utterance}. In addition, the interaction module maintains child-specific context through an RDF/SPARQL memory for preferences and a JSON session-score history. The RDF memory is used to personalize the child-facing dialogue, while session scores are updated after the activity and are later queried by the orchestrator to propose the activity that previously worked best for the same child.

Dynamic updates are handled at mission boundaries and at interaction events. The static scene graph is intentionally stable: it is a prior for a semi-structured therapy room, and changes such as moving the therapist or child location are made by editing the corresponding graph nodes rather than by continuous tracking. The dynamic visual graph, instead, is regenerated whenever the EMPOWER capture/planning pipeline is run. Therefore, object relocation is handled by re-scanning the scene: new detections replace the previous \texttt{global\_scene\_graph.json}, moved objects receive updated centroids, and newly visible objects can be added to the current grounding state. During execution, simulated manipulation also changes the Gazebo world state by deleting a grabbed object and respawning it at the resolved drop pose, while the executor tracks the currently held object internally. New user intentions are handled online by the mission state machine: a therapist override replaces the proposed activity/task before planning, the updated task description is sent to EMPOWER, and a new symbolic plan is generated from the current instruction and current scene grounding. The system does not implement persistent continuous semantic tracking of all objects over long time horizons; instead, it uses the static graph for durable social/task priors, the regenerated EMPOWER graph for the current visual state, and interaction memories for longitudinal personalization across therapy sessions.
